\label{sec:Project-Details}
The simulations of spatio-temporal dynamics outlined in Sec.~\ref{sec:research_plan} will be carried out with the multiphysics simulation package uG4 \citep{Vogel2013, Heppner13}. This decades-old project focusses on solving partial differential equations on unstructured grids and has been successfully used in applications from diffusion-advection systems, Navier-Stokes equations, and mechanics applications. In particular the code has been developed to run very well on highly parallel architectures \citep{Vogel2013, Heppner13}. To document uG4's general scaling behavior, we performed benchmark scaling simulations to solve the second order elliptic partial differential equation $\Delta \Phi = 0$, the Laplace equation with Dirichlet boundary conditions. The PDE problem was solved on the unit cube with seven levels of regular grid refinement which contains 16,974,593 DoFs. The scaling is shown in Table ~\ref{tab:xsede_strong_scaling} of \textbf{Section General scaling properties of uG4} in the \textit{performance} document. For this proposal we will be making use of first-order finite volume discretization methods for the underlying systems of partial differential equations. The resulting system of linear equations will be solved using a geometric multigrid method (GMG), combined with an implicit time-stepping scheme. For more details we refer to Sec.~\ref{sec:research_plan}.
%The simulation framework uG4 is publicly available on Github\footnote{See the various repositories in the organization \texttt{UG4} on \url{https://github.com/UG4}}. Upon peer-reviewed publication all relevant scripts for data reproduction will be made publicly available through a repository in the existing Github organization.
%Produced simulation results are stored in VTK format (1 file per time step). After completion of a simulation run, e.g. parameter sweeps, we will copy data to a local storage facility. Over the period of an average month we expect to accumulate $<$ 1 TiB of VTK data. 
